\chapter{Formal Language Theory of Transformers}
\label{ch11:top}

This chapter has one guiding question: what do transformer expressivity results suggest about the kinds of generating functions one should expect from transformer-based language models?

The word \emph{suggest} matters. Results in formal-language theory and circuit complexity are primarily about \emph{recognition} and \emph{simulation}: what classes of languages a model can recognize, or what computational devices it can simulate. Generating-function claims are about \emph{counting sequences} or \emph{output distributions}. The bridge between these two worlds is often illuminating, but it is not automatic. So this chapter is best read as a survey plus a modeling stance, not as a chapter that fully proves the analytic conclusions it points toward.

\section{Circuit Classes as a Measuring Stick}

The complexity classes that appear most often in transformer expressivity papers are the following.

\begin{definition}
\begin{itemize}
  \item $\mathsf{AC}^0$: constant-depth, polynomial-size Boolean circuits with unbounded-fan-in AND, OR, and NOT gates.
  \item $\mathsf{TC}^0$: the same, but with threshold gates added.
  \item $\mathsf{NC}^1$: logarithmic-depth, polynomial-size circuits with bounded fan-in.
\end{itemize}
\end{definition}

Intuitively, $\mathsf{AC}^0$ can aggregate many bits quickly but is weak at global counting. A classical example is parity: deciding whether a binary string has an even number of ones is not in $\mathsf{AC}^0$, but it is in $\mathsf{TC}^0$ by a constant-depth threshold-circuit construction. It is \emph{not} computed by a single threshold gate; the point is only that threshold gates make parity easy in constant depth.

The class $\mathsf{NC}^1$ is stronger because logarithmic depth allows recursive composition. Many expression-evaluation and parsing tasks naturally live there. In this chapter we suppress uniformity details; the cited papers make those conventions precise.

\section{Single-Pass Transformers and the $\mathsf{TC}^0$ Heuristic}

A strong line of work, surveyed in \cite{StroblEtAl2024}, shows that several fixed-depth transformer variants sit inside or near $\mathsf{TC}^0$-like expressivity classes when one studies single-pass recognition.

\begin{theorem}[Merrill, Sabharwal, Smith 2022, black box]
Saturated transformers of constant depth can be simulated by polynomial-size constant-depth threshold circuits. In particular, the languages they recognize lie in $\mathsf{TC}^0$.
\end{theorem}

The point for us is not the proof technology. It is the resulting modeling picture: fixed-depth single-pass transformers appear much closer to finite-state or threshold-circuit behavior than to full context-free parsing power.

But this is only a starting point for analytic combinatorics. A recognition upper bound does \emph{not} by itself prove that an output distribution has a rational generating function. What it does is motivate a working hypothesis:
\begin{quote}
For single-pass transformers, finite-state / WFA-like approximations are plausible enough to study seriously.
\end{quote}
That is a modeling stance, not a theorem yet.

\section{Impossibility Results as Warnings}

Hahn's theorem gives a complementary message.

\begin{theorem}[Hahn 2020, black box]
Fixed-depth soft-attention transformers cannot uniformly solve certain length-unbounded tasks such as parity and multi-type Dyck recognition.
\end{theorem}

The safest interpretation is architectural, not analytic. These theorems say that some globally sensitive or deeply nested behaviors are hard for bounded-depth soft-attention models to represent uniformly across all lengths. They do \emph{not} directly tell us that a particular pole or branch point must or must not appear in a generating function.

So the correct lesson is modest: lower bounds on recognition complexity warn us against treating single-pass transformers as arbitrary pushdown-like devices. They do not by themselves determine the analytic type of an output distribution.

\section[Exact Rational Regimes]{Exact Rational Regimes: RNNs and Automata}

There are, however, exact positive results on the neural side.

\begin{theorem}[Svete and Cotterell 2023, black box]
A certain class of recurrent language models is exactly equivalent to a corresponding class of probabilistic finite-state automata.
\end{theorem}

This matters because Chapter~7 already told us what finite-state weighted models imply analytically: rational length generating functions, pole-based asymptotics, and a finite-state matrix description. So the theorem shows that at least some neural sequence models do sit exactly inside the rational regime.

That does \emph{not} mean all recurrent models do, and it certainly does not mean all transformer models do. The safe conclusion is narrower: there is an exact neural-to-rational corner of the landscape, so rational approximations are not an empty fantasy.

\section{Transformers Simulating Weighted Automata}

Another important direction is constructive simulation.

\begin{theorem}[Rizvi et al. 2024, black box]
Transformers with suitable attention mechanisms can exactly or approximately simulate weighted finite automata on inputs up to a prescribed length bound.
\end{theorem}

Again, the correct analytic takeaway is limited. Exact simulation of a WFA means rational behavior for that simulated system. Approximate simulation on strings up to length $T$ means that finite-state behavior can be reproduced well on a finite range. It does \emph{not} by itself imply that the transformer's global generating function has approximately the same singularities or asymptotics as the WFA for all lengths.

So the right lesson is:
\begin{itemize}
  \item exact WFA behavior is available to transformers,
  \item finite-length WFA approximation is a meaningful program,
  \item but asymptotic analytic type is a stronger question than finite-length simulation.
\end{itemize}

That stronger question will be taken up in Chapter~12 rather than assumed here.

\section{Chain of Thought Changes the Picture}

All of the previous discussion concerns models that map an input to an output in a bounded-depth forward pass. Chain-of-thought changes that regime by allowing the model to emit intermediate tokens and then condition on them.

\begin{theorem}[Merrill and Sabharwal 2024, black box]
Under the hypotheses of their framework, a transformer with polynomially many chain-of-thought steps can realize computations all the way up to polynomial time.
\end{theorem}

For our purposes, the crucial implication is qualitative. Once intermediate computation is allowed, the finite-state / rational heuristic becomes much less compelling. The model may now simulate much richer mechanisms than a single-pass bounded-depth device.

What should one conclude analytically? Not that chain-of-thought automatically implies algebraic generating functions or an $n^{-3/2}$ law. That would be far too strong. The safe statement is only this:
\begin{quote}
Chain-of-thought makes context-free-like, pushdown-like, and even richer behavior much more plausible than in the single-pass regime.
\end{quote}
Whether that plausibility shows up in the actual output-support or output-probability generating functions of a trained model is a separate question.

\section{A Three-Level Synthesis}

The chapter's message can be compressed into three levels of confidence.

\begin{enumerate}
  \item \textbf{Theorems.} We have cited exact results about recognition bounds, exact finite-state equivalences, and exact or finite-length simulation theorems.
  \item \textbf{Modeling stance.} For single-pass transformers, it is reasonable to investigate rational / WFA-style approximations first.
  \item \textbf{Open analytic question.} How those expressivity facts constrain the actual generating functions of output distributions is not settled by this chapter alone.
\end{enumerate}

This is the stance the rest of the book will adopt. Single-pass models are treated as candidates for rational approximation. Chain-of-thought models are treated as candidates for richer algebraic or beyond-algebraic behavior. Both are hypotheses to be tested, not conclusions already extracted from the cited complexity theorems.

\section*{Preview}

Chapter~12 turns this survey perspective into an operational question: if a transformer is WFA-like in a useful sense, what exactly should be approximated, and what analytic conclusions survive that approximation? That is where the approximation target must finally be made precise.

\section*{Further Reading}

\begin{itemize}
  \item \textbf{L. Strobl et al.}, ``What formal languages can transformers express? A survey'' (2024) \cite{StroblEtAl2024} --- the best overview of transformer expressivity results.
  \item \textbf{W. Merrill, A. Sabharwal, and N. A. Smith}, ``Saturated transformers are constant-depth threshold circuits'' (2022) \cite{MerrillSabharwalSmith2022} --- the foundational $\mathsf{TC}^0$ upper-bound paper.
  \item \textbf{M. Hahn}, ``Theoretical limitations of self-attention in neural sequence models'' (2020) \cite{Hahn2020} --- for parity and Dyck-style lower bounds on soft-attention models.
  \item \textbf{A. Svete and R. Cotterell}, ``Recurrent neural networks and weighted finite-state transducers'' (2023) \cite{SveteCotterell2023} --- for an exact finite-state corner of neural sequence modeling.
  \item \textbf{M. Rizvi-Martel et al.}, ``Simulating weighted automata over sequences and trees with transformers'' (2024) \cite{RizviEtAl2024} --- for constructive transformer simulation of weighted automata.
  \item \textbf{W. Merrill and A. Sabharwal}, ``The expressive power of transformers with chain of thought'' (2024) \cite{MerrillSabharwal2024} --- for the chain-of-thought expressivity jump.
\end{itemize}
